\chapter{Instructions to Run System}
\label{apx:instructions}

\section{Environment Setup}

The system requires a VLLM server to serve the VLM. It uses a Conda environment with specific packages. The environment configuration is provided as a YAML file (\texttt{environment\_vllm.yml}).
The VLLm server was run on an A100 40GB GPU with driver version - and cuda version - 13.0

\subsection{Prerequisites}

\begin{itemize}
    \item Conda (Miniconda or Anaconda) installed on the system.
    \item NVIDIA GPU with CUDA support (A100 40GB or 80GB recommended).
    \item Linux operating system (tested on Ubuntu).
\end{itemize}

\subsection{Installation Steps}

\begin{enumerate}
    \item \textbf{Clone the repository:}
\begin{verbatim}
git clone <repository-url>
cd Project
\end{verbatim}

    \item \textbf{Create the Conda environment from the YAML file:}
\begin{verbatim}
conda env create -f environment_vllm.yml
\end{verbatim}

    \item \textbf{Activate the environment:}
\begin{verbatim}
conda activate vllm
\end{verbatim}

    \item \textbf{Verify the installation:}
\begin{verbatim}
python -c "import vllm; print(vllm.__version__)"
python -c "import torch; print(torch.cuda.is_available())"
\end{verbatim}
\end{enumerate}

\subsection{Alternative Installation (from spec file)}

If the YAML installation fails due to platform differences, use the spec file:

\begin{verbatim}
conda create --name vllm --file conda_vllm_packages.txt
\end{verbatim}

\subsection{Running the VLLM server}

After activating the environment, the vllm server can be run using:

\begin{verbatim}
srun -p gpu -w gorina8 --gres=gpu:1 --pty $SHELL
\end{verbatim}

and after getting access to the server, run the following commands

\begin{verbatim}
conda activate vllm

vllm serve Qwen/Qwen3-VL-8B-Instruct \
  --tensor-parallel-size 1 \
  --max-model-len 32768 \
  --gpu-memory-utilization 0.90 \
  --limit-mm-per-prompt '{"video":0, "image":10}' \
  --max-num-seqs 128 \
  --port 8005
\end{verbatim}

\subsection{Running the pipeline using Iterative Question-Answer generation}

For running the pipeline using the Ground Truth questions, we need to set GT\_EVID=True and GT\_QUES=True in \begin{verbatim}run_with_vllm_splade_bge_longclip_icl_qwen_iterative.sh\end{verbatim}

and get access to a gpu server using
\begin{verbatim}
srun -p gpu -w gorina8 --gres=gpu:1 --pty $SHELL
\end{verbatim}

and run 

\begin{verbatim}
run_with_vllm_splade_bge_longclip_icl_qwen_iterative.sh
\end{verbatim}

If we want the VLM to generate the questions we can set GT\_EVID=False and GT\_QUES=False

\subsection{Running the pipeline using claim based retrieval in Question-Answer generation}

Get access to a gpu server using
\begin{verbatim}
srun -p gpu -w gorina8 --gres=gpu:1 --pty $SHELL
\end{verbatim}

and run 

\begin{verbatim}
run_with_vllm_splade_bge_longclip_icl_qwen.sh
\end{verbatim}
